\chapter{The classical mod--2 Adams spectral sequence}
\label{chap:classical-adams}

This chapter turns the paper's classical background into explicit interfaces.
The high-stem Ext calculations remain evidence; only general algebra and the
paper's internal deductions are project proof obligations.

\section{Steenrod algebra, Ext, and convergence}

% SOURCE: aimpaper/main.tex:126-150; PROJECT_BOUNDARY.md, Steenrod algebra and Ext.
\begin{definition}[Dual Steenrod algebra]
  \label{def:dual-steenrod-algebra}
  \uses{def:mathlib-f2}
  \notready
  The project provides the graded connected Hopf algebra
  $\mathcal A_*=H_*(H\mathbb F_2;\mathbb F_2)$, its augmentation and coproduct,
  and the graded-dual relation with the mod--$2$ Steenrod algebra $\mathcal A$.
  Mathlib v4.32.2 has no Steenrod-algebra implementation that supplies this
  structure.
\end{definition}

% SOURCE: aimpaper/main.tex:126-150 and every Adams E_2-page computation.
\begin{definition}[Graded Steenrod-comodule category]
  \label{def:steenrod-comodule-category}
  \uses{def:dual-steenrod-algebra,def:mathlib-abelian,def:hf2-homology}
  \notready
  Graded right $\mathcal A_*$-comodules form the abelian category used to model
  mod--$2$ Adams input.  The interface includes degree-preserving comodule
  maps, shifts, kernels, cokernels, tensor products, the trivial comodule
  $\mathbb F_2$, and the homology comodule $H_*(X;\mathbb F_2)$.
\end{definition}

% SOURCE: aimpaper/main.tex:126-150; internal and cohomological Ext gradings.
\begin{definition}[Bigraded Adams Ext]
  \label{def:adams-bigraded-ext}
  \uses{def:steenrod-comodule-category,def:mathlib-derived-ext}
  \notready
  The Adams group $\operatorname{Ext}_{\mathcal A}^{s,t}(M,\mathbb F_2)$ has
  derived degree $s$ and internal degree $t$.  It is represented equivalently
  by internal graded Ext in $\mathcal A_*$-comodules, with an explicit proof
  translating the paper's module convention to the comodule convention.  The
  interface supplies Yoneda products and internal-degree shifts; Mathlib's
  singly graded derived Ext alone does not.
\end{definition}

% SOURCE: aimpaper/main.tex:126-150; finite-type standing assumption at lines 253-263.
\begin{proposition}[Module--comodule Adams Ext comparison]
  \label{prop:adams-module-comodule-ext-comparison}
  \uses{def:adams-bigraded-ext,def:admissible-spectrum,def:hf2-homology}
  \notready
  For an admissible spectrum $X$, degreewise finite duality identifies the
  paper's module expression
  \[
    \operatorname{Ext}_{\mathcal A}^{s,t}
      (H^*(X;\mathbb F_2),\mathbb F_2)
  \]
  with the corresponding internal Ext, equivalently Cotor, expression built
  from the graded $\mathcal A_*$-comodule $H_*(X;\mathbb F_2)$.  The
  isomorphism preserves $(s,t)$, products, suspension, and maps induced by
  spectra.
\end{proposition}

\begin{proof}
  Use finite type to dualize each graded piece, identify Steenrod-module
  actions with dual-Steenrod coactions, and compare the bar and cobar
  resolutions degreewise.  Verify the internal grading and Yoneda product on
  the chain-level comparison before passing to cohomology.
\end{proof}

% SOURCE: aimpaper/main.tex:134-150; PROJECT_BOUNDARY.md, Steenrod algebra and Ext.
\begin{definition}[Mod--2 Steenrod algebra and Ext]
  \label{def:steenrod-ext}
  \uses{def:adams-bigraded-ext,prop:adams-module-comodule-ext-comparison,def:hf2-homology}
  \notready
  Let $\mathcal A$ be the mod--$2$ Steenrod algebra and let
  $\mathbb F_2$ be its trivial module.  The project interface provides the
  bigraded groups
  $\operatorname{Ext}_{\mathcal A}^{s,t}(H^*(X;\mathbb F_2),\mathbb F_2)$,
  Yoneda products, suspension/regrading maps, and functoriality in $X$.
\end{definition}

% SOURCE: aimpaper/main.tex:126-150 and 253-275; standard mod-2 Adams resolution used implicitly.
\begin{definition}[$H\mathbb F_2$-Adams resolution and tower]
  \label{def:hf2-adams-tower}
  \uses{def:admissible-spectrum,def:hf2-homology,def:cofiber-triangle,def:stable-sphere-suspension}
  \notready
  For an admissible spectrum $X$, iterate the unit
  $X\to H\mathbb F_2\wedge X$ and its fiber to obtain a functorial Adams
  resolution and a decreasing tower
  \[
    \cdots\longrightarrow X_2\longrightarrow X_1\longrightarrow X_0=X.
  \]
  The interface records the exact triangles, augmentations to $X$, naturality,
  multiplicative pairings, and the fact that every successive cofiber is an
  $H\mathbb F_2$-module spectrum.
\end{definition}

% SOURCE: aimpaper/main.tex:253-343, Adams filtration on classes and maps.
\begin{definition}[Filtration induced by the Adams tower]
  \label{def:adams-tower-filtration}
  \uses{def:hf2-adams-tower,def:stable-homotopy-group,def:core-filtration}
  \notready
  The Adams filtration on $\pi_nX$ is
  \[
    F^s\pi_nX=\operatorname{im}\bigl(\pi_nX_s\longrightarrow\pi_nX\bigr).
  \]
  A map has filtration at least $k$ when it factors through the $k$th stage of
  the target tower, with the chosen factorization retained as data.  Tower
  pairings give the filtration inequalities for products and composition.
\end{definition}

% SOURCE: aimpaper/main.tex:126-150 and standing convergence assumption at 253-263; classical Adams construction.
\begin{theorem}[Adams tower, Ext, and convergence comparison]
  \label{thm:adams-tower-ext-comparison}
  \uses{def:hf2-adams-tower,def:adams-tower-filtration,def:steenrod-ext,thm:filtered-complex-to-spectral-sequence,def:strong-convergence-detection,def:external-result}
  \notready
  The exact couple or filtered cobar complex of the
  $H\mathbb F_2$-Adams tower yields a multiplicative spectral sequence whose
  $E_2$ page is naturally
  $\operatorname{Ext}_{\mathcal A}^{s,t}
  (H^*(X;\mathbb F_2),\mathbb F_2)$ and whose induced abutment filtration is
  Definition~\ref{def:adams-tower-filtration}.  A located classical Adams
  convergence result supplies strong convergence for the admissible spectra
  used in the paper.
\end{theorem}

\begin{proof}
  Apply the filtered-complex construction to the tower's cobar resolution,
  identify its normalized cochains with the Steenrod cobar complex, and use
  the module--comodule comparison to obtain the bigraded Ext page.  Transport
  tower pairings to page products.  Finally instantiate the located
  convergence theorem and identify its complete separated filtration with the
  tower-image filtration.
\end{proof}

% SOURCE: aimpaper/main.tex:134-150 and 253-263.
\begin{definition}[Classical Adams spectral sequence]
  \label{def:classical-adams-ss}
  \uses{thm:adams-tower-ext-comparison,def:steenrod-ext,def:adams-bidegree,def:adams-complex-shape,def:chosen-page-presentation,def:multiplicative-ss,def:admissible-spectrum,def:stable-homotopy-group,def:strong-convergence-detection}
  \notready
  For every admissible spectrum $X$, the classical mod--$2$ Adams spectral
  sequence has
  \[
    E_2^{s,t}(X)=\operatorname{Ext}_{\mathcal A}^{s,t}
      (H^*(X;\mathbb F_2),\mathbb F_2),
    \qquad d_r:E_r^{s,t}\to E_r^{s+r,t+r-1},
  \]
  together with the chosen $Z_r/B_r$ presentation, products, naturality, and
  strong convergence to the $2$-complete stable homotopy of $X$.
\end{definition}

% SOURCE: aimpaper/main.tex:253-263 and 841-864.
\begin{proposition}[Classical Adams convergence and detection]
  \label{prop:classical-adams-convergence}
  \uses{def:classical-adams-ss,def:strong-convergence-detection,def:adams-detection}
  \notready
  The $E_\infty$ page of the classical Adams spectral sequence is naturally
  the associated graded of the Adams filtration on the $2$-complete stable
  homotopy groups.  Consequently its permanent representatives detect exactly
  the corresponding filtered homotopy classes.
\end{proposition}

\begin{proof}
  Apply the strong-convergence component of the classical Adams interface and
  identify its filtration with the Adams-tower filtration.  The chosen
  $Z_\infty/B_\infty$ presentation then gives the stated detection relation
  degree by degree.
\end{proof}

% SOURCE: aimpaper/main.tex:136-162.
\begin{definition}[The classes $h_j$]
  \label{def:adams-hj}
  \uses{def:classical-adams-ss}
  \notready
  For $j\ge0$, $h_j$ denotes the standard generator in
  $E_2^{1,2^j}(S^0)$.  Its square $h_j^2$ lies in
  $E_2^{2,2^{j+1}}(S^0)$ and has stem $2^{j+1}-2$.  In particular,
  $h_6^2$ lies in bidegree $(2,128)$ and stem $126$.
\end{definition}

\section{Low-line literature inputs}

% SOURCE: aimpaper/main.tex:140-150; citations Adams/Mahowald--Tangora/May.
% VERIFIED-IN: aimpaper/main.tex:140-150 contains a typo: survival is j <= 3, while d_2 is nonzero for j >= 4.
\begin{theorem}[Adams one-line differentials]
  \label{thm:external-adams-one-line}
  \uses{def:adams-hj,def:external-result}
  \notready
  An explicit external result supplies that $h_0,h_1,h_2,h_3$ survive and
  that, for $j\ge4$,
  \[
    d_2(h_j)=h_0h_{j-1}^2\ne0.
  \]
  This is the mathematically consistent reading of the two consecutive
  sentences in \texttt{aimpaper/main.tex:140--150}; the printed ``$j\ge3$
  survives'' is recorded as a source typo, not adopted as a project theorem.
\end{theorem}

\begin{proof}
  This node is a literature input.  Instantiate the external-result record at
  the cited theorem and transport its grading to
  Definition~\ref{def:adams-bidegree}; no internal proof of the classical
  calculation is claimed.
\end{proof}

% SOURCE: aimpaper/main.tex:151-163, Corollary cor:2line and its cited inputs.
\begin{corollary}[Survivors on the Adams two-line]
  \label{cor:adams-two-line-survivors}
  \uses{def:adams-hj,thm:external-adams-one-line,thm:h6-square-permanent,def:external-result,def:external-evidence}
  \notready
  Subject to the explicitly cited Ext calculations and differential results,
  the nonzero classes on the Adams two-line that survive to $E_\infty$ are
  precisely
  \[
    h_0h_2,\quad h_0h_3,\quad h_2h_4,\quad
    h_1h_j\ (j\ge3),\quad h_j^2\ (0\le j\le6).
  \]
\end{corollary}

\begin{proof}
  The supplied Adams, Lin, May, Mahowald--Tangora, Isaksen--Wang--Xu,
  Mahowald, Barratt--Jones--Mahowald, and Hill--Hopkins--Ravenel results first
  enumerate the $E_3$ candidates and then remove $h_2h_5$, $h_3h_5$,
  $h_3h_6$, and every $h_j^2$ with $j\ge7$.  Their positive survival results
  give every listed class except $h_6^2$; apply
  Theorem~\ref{thm:h6-square-permanent} for that final class.  The dependency
  points from this corollary to the main theorem, never conversely.
\end{proof}

% SOURCE: aimpaper/main.tex:164-171, Questions que:2theta6 and que:theta5sq.
\begin{definition}[Order-two $\theta_6$ question]
  \label{def:question-order-two-theta6}
  \uses{def:open-proposition,def:adams-hj,def:adams-detection}
  \notready
  This open proposition asks whether $h_6^2$ detects a Kervaire class
  $\theta_6\in\pi_{126}S^0$ of exact order two.  It is recorded but is not a
  hypothesis of the permanent-cycle proof.
\end{definition}

% SOURCE: aimpaper/main.tex:168-171, Question que:theta5sq.
\begin{definition}[$\theta_5^2=0$ question]
  \label{def:question-theta5-square-zero}
  \uses{def:open-proposition,def:adams-detection}
  \notready
  This open proposition asks whether one can choose a class
  $\theta_5\in\pi_{62}S^0$ detected by $h_5^2$ such that
  $\theta_5^2=0$.  It is neither assumed nor proved by this project.
\end{definition}
